\providecommand{\kBKM}{\kappa_{\mathrm{BKM}}}
\providecommand{\kch}{\kappa_{\mathrm{ch}}}
\providecommand{\kcat}{\kappa_{\mathrm{cat}}}
\providecommand{\kfib}{\kappa_{\mathrm{fiber}}}
\providecommand{\Z}{\mathbb{Z}}
\providecommand{\C}{\mathbb{C}}
\providecommand{\DT}{\mathrm{DT}}
\providecommand{\DTred}{Z^{\mathrm{red}}_{\mathrm{DT}}}
\providecommand{\PTred}{Z^{\mathrm{red}}_{\mathrm{PT}}}
\providecommand{\GWred}{Z^{\mathrm{red}}_{\mathrm{GW}}}
\providecommand{\Dfive}{D_5}
\providecommand{\DeltaFive}{\Delta_5}
\providecommand{\PhiTenUn}{\Phi_{10}^{\mathrm{un}}}
\providecommand{\PhiTenOP}{\Phi_{10}^{\mathrm{OP}}}
\providecommand{\ZOP}{Z_{\mathrm{OP}}}
\providecommand{\chiTen}{\chi_{10}}
\providecommand{\phizo}{\phi_{0,1}}
\providecommand{\BKM}{\mathrm{BKM}}
\providecommand{\CoHA}{\mathrm{CoHA}}
\providecommand{\Yplus}{Y^+}
\providecommand{\Eone}{E_1}
\providecommand{\Etwo}{E_2}
\providecommand{\Woneinfty}{\mathcal W_{1+\infty}}

\section*{Oberdieck-family anchors for \(K3\times E\) and CHL reduced DT}
\label{sec:oberdieck-family-primary-literature}

Let \(S\) be a projective K3 surface, let \(E\) be an elliptic curve, and
put \(X=S\times E\). The enumerative input has four independent layers:
the primitive Igusa theorem for \(S\times E\), the
Oberdieck-Pandharipande monic reduced-DT scalar convention, the CHL
primitive scalar conjecture with its boundary coefficients, and the
imprimitive multiple-cover problem. The scalar normalisation is
\[
  \Dfive = 64^{-1}\DeltaFive,\qquad
  \PhiTenUn = \DeltaFive^2,\qquad
  \PhiTenOP = \Dfive^2 = 4096^{-1}\DeltaFive^2 .
\]
The Behrend-weighted reduced Donaldson-Thomas scalar in this monic
convention is
\[
  \ZOP
  =
  -(\PhiTenOP)^{-1}
  =
  -\Dfive^{-2}
  =
  -4096\,\DeltaFive^{-2}.
\]
Thus \(\PhiTenUn=\DeltaFive^2\) is the unnormalised square, while
\(\PhiTenOP=4096^{-1}\DeltaFive^2\) is the scalar used by the
Oberdieck-Pandharipande convention. The primitive \(\BKM\) denominator is
the square root \(\DeltaFive\):
\[
  \kBKM(\DeltaFive)
  =
  \frac{c_1(0)}{2}
  =
  \frac{10}{2}
  =
  5 .
\]
The compact \(K3\times E\) entries are
\[
  \kcat(K3\times E)=0,\qquad
  \kch(K3\times E)=0,\qquad
  \kch^{\mathrm{Heis}}(K3\times E)=3,\qquad
  \kBKM(\DeltaFive)=5,\qquad
  \kfib(K3\times E)=24 .
\]
The total-space categorical value is \(0\); the value
\(2=\chi(\mathcal O_{K3})\) is solely the K3-fibre value. These five
numbers are independent readings: Hodge supertrace, compact chiral
supertrace, Heisenberg-Mukai specialisation, Borcherds weight, and
Mukai-lattice fibre rank.

\subsection*{Primitive \(K3\times E\) and the Igusa scalar}

\emph{Primary source.} G.~Oberdieck and A.~Pixton,
\emph{Holomorphic anomaly equations and the Igusa cusp form
conjecture}, \emph{Invent.\ Math.} \textbf{213} (2018), 507-587,
arXiv:1706.10100.

\emph{Theorem (Oberdieck-Pixton, Thm.~1).}
Let \(\beta_h\in H_2(S,\Z)\) be primitive with
\(\langle\beta_h,\beta_h\rangle=2h-2\). For the reduced
Gromov-Witten invariants \(N_{g,h,d}\) of \(S\times E\), the
primitive partition function
\[
  Z(u,q,\widetilde q)
  =
  \sum_{g,h,d} N_{g,h,d}\,
  u^{2g-2}q^{h-1}\widetilde q^{d-1}
\]
is the Laurent expansion
\[
  Z(u,q,\widetilde q)
  =
  -\chiTen(p,q,\widetilde q)^{-1},
  \qquad p=e^{iu}.
\]
The cusp form \(\chiTen\) is the genus-two weight-ten Igusa form in the
product normalisation used by the \(K3\times E\) curve-counting
theorem. With the scalar notation above,
\[
  \chiTen=\PhiTenOP=\Dfive^2.
\]
Consequently the reduced scalar attached to the primitive \(K3\)-class
sector is the OP-normalised inverse square
\[
  -\chiTen^{-1}
  =
  -(\PhiTenOP)^{-1}
  =
  -4096\,\DeltaFive^{-2}.
\]

\emph{Boundary.} The theorem is primitive in the K3 factor. It proves a
scalar GW/PT Igusa identity. The BKM object extracted from the
Borcherds denominator is the weight-five datum \(\DeltaFive\), while
the OP scalar is its monic square.

\subsection*{Primitive BKM denominator rows}

\emph{Primary sources.} R.~E.~Borcherds,
\emph{Automorphic forms with singularities on Grassmannians},
\emph{Invent.\ Math.} \textbf{132} (1998), 491-562, Thm.~13.3;
V.~Gritsenko and V.~Nikulin, \emph{Siegel automorphic form corrections
of some Lorentzian Kac-Moody Lie algebras}, \emph{Amer.\ J.\ Math.}
\textbf{119} (1997), 181-224; V.~Gritsenko, \emph{Elliptic genus of
Calabi-Yau manifolds and Jacobi and Siegel modular forms}, \emph{St.\
Petersburg Math.\ J.} \textbf{11} (2000), 781-804.

Borcherds' singular theta correspondence gives the weight of an
automorphic product as one half of the zeroth Fourier coefficient of
the input at the cusp:
\[
  \operatorname{wt}(\Phi_N)=\frac{c_N(0)}{2}.
\]
For the primitive compact CHL denominator ladder the row label is
\[
  N\in\{1,2,3,4,6\},
\]
and the Gritsenko-Nikulin-Gritsenko products give
\[
\begin{array}{c|c|c|c}
N & \text{denominator} & c_N(0) & \kBKM \\ \hline
1 & \Delta_5=\Phi_1 & 10 & 5\\
2 & \Delta_4^{(2)} & 8 & 4\\
3 & \Delta_3^{(3)} & 6 & 3\\
4 & \Delta_2^{(4)} & 4 & 2\\
6 & \Delta_1^{(6)} & 2 & 1
\end{array}
\]
Equivalently,
\[
  (c_N(0))_{N=(1,2,3,4,6)}=(10,8,6,4,2),\qquad
  (\kBKM(\Phi_N))_{N=(1,2,3,4,6)}=(5,4,3,2,1).
\]
The orders \(N=5,7,8\) require a separate theorem supplying a common
compact host and a map to the five-row denominator ladder.

\emph{Catalogue boundary.} The primitive CHL denominator ladder, the
Gritsenko-Clery diagonal-divisor catalogue, the scalar inverse-square
convention, and compact Hall or \(\BKM\) recognition are separate data.
The Gritsenko-Clery catalogue is triple-indexed by \((t,N;k)\), has
weight sequence
\[
  (5,2,3,1,2,\tfrac12,\tfrac32,1),
\]
and has twice-weight sequence
\[
  (10,4,6,2,4,1,3,2).
\]
Its common entry with the primitive CHL ladder is the level-one
\(\Delta_5\) row.

\subsection*{CHL primitive scalar and boundary coefficients}

\emph{Primary source.} J.~Bryan and G.~Oberdieck,
\emph{CHL Calabi-Yau threefolds: curve counting, Mathieu moonshine
and Siegel modular forms}, \emph{Commun.\ Number Theory Phys.}
\textbf{14} (2020), 785-862, arXiv:1811.06102.

For an elliptic CHL model \(X_N=[S\times E/\langle(g_N,t_N)\rangle]\),
Bryan-Oberdieck define the primitive reduced Donaldson-Thomas series
\[
  Z^{X_N}(q,t,p)
  =
  \sum_{h,d,n}
  \DT^{X_N}_{n,(\beta_h,d)}\,
  q^{d-1}t^{\langle\beta_h,\beta_h\rangle/2}(-p)^n .
\]
Let \(\widetilde\Phi_N\) be the Borcherds lift of the corresponding
twisted-twined K3 elliptic genera.

\emph{Conjecture A (Bryan-Oberdieck).}
Under \(q=e^{2\pi i\tau}\),
\(t=e^{2\pi i\sigma}\), \(p=e^{2\pi iz}\),
\[
  Z^{X_N}(q,t,p)
  =
  -\widetilde\Phi_N(\tau,z,\sigma)^{-1}.
\]
The scalar \(\widetilde\Phi_N\) is a Siegel modular form of weight
\[
  \left\lceil \frac{24}{N+1}\right\rceil -2
\]
in the Bryan-Oberdieck CHL convention. At \(N=1\) this specialises to
the Igusa form \(\chiTen=\PhiTenOP\).

\emph{Theorem 0.1 (Bryan-Oberdieck).}
For an elliptic CHL model of order \(N\), the two
boundary coefficients of \(Z^{X_N}\) satisfy
\[
  [Z^{X_N}(q,t,p)]_{t^{-1/N}}
  =
  \frac{1}{\Theta(q^N,p)^2\Delta_N(q)},
  \qquad
  [Z^{X_N}(q,t,p)]_{q^{-1}}
  =
  \frac{1}{\Theta(t,p)^2\Delta_N(t^{1/N})}.
\]
Thus Conjecture A is proved after extracting the first coefficient in
the \(t\)-direction and the first coefficient in the \(q\)-direction.
The proof uses the orbifold topological vertex for the \(t^{-1/N}\)
coefficient and a degeneration to \(K3\times\mathbb P^1\) for the
\(q^{-1}\) coefficient.

\emph{Boundary.} Theorem~0.1 is a boundary-coefficient theorem for the
scalar DT series. The all-coefficient primitive series for \(N\ge2\), higher
primitive squares, and imprimitive classes remain governed by
Conjecture A and the multiple-cover problem.

\subsection*{CHL dyon forms and Mathieu twining}

\emph{Primary source.} A.~Chattopadhyaya and J.~R.~David,
\emph{Dyon degeneracies from Mathieu moonshine symmetry},
\emph{Phys.\ Rev.\ D} \textbf{96} (2017), 086020, arXiv:1704.00434.

\emph{Chattopadhyaya-David.}
Chattopadhyaya and David construct the twisted elliptic genera in all
sectors for the relevant
Mathieu classes, form the corresponding theta-lift Siegel modular
forms, and verify the predicted integrality and sign properties of the
inverse Fourier expansions for \(1/4\)-BPS dyons in type-II
compactifications on \((K3\times T^2)/\Z_N\). The result supplies the
seven CHL dyon partition functions and the Mathieu-twined input used by
the CHL scalar forms.

\emph{Boundary.} The result is a string-theoretic BPS-index statement.
It checks the automorphic scalar forms and their Fourier coefficients;
the reduced DT theorem for \(X_N\) and compact Hall/\(\BKM\) generators
and relations are separate mathematical obligations. The dyon rows
\(N=5,7,8\) stay in this physics row system until a compact CHL
denominator theorem matches them to the five-row ladder above.

\subsection*{K3 BPS counts and the KKV theorem}

\emph{Primary source.} R.~Pandharipande and R.~P.~Thomas,
\emph{The Katz-Klemm-Vafa conjecture for K3 surfaces},
\emph{Forum Math.\ Pi} \textbf{4} (2016), e4, arXiv:1404.6698.

\emph{Theorem 1 (Pandharipande-Thomas).}
For a K3 surface \(S\), the BPS count
\(r_{g,\beta}\) depends only on the norm
\(\langle\beta,\beta\rangle=2h-2\), and the Katz-Klemm-Vafa formula
holds:
\[
  \sum_{g\ge0}\sum_{h\ge0}
  (-1)^g r_{g,h}
  \bigl(y^{1/2}-y^{-1/2}\bigr)^{2g}q^h
  =
  \prod_{n\ge1}
  \frac{1}{(1-q^n)^{20}(1-yq^n)^2(1-y^{-1}q^n)^2}.
\]

\emph{Boundary.} This is a theorem for reduced K3 curve counts, including
imprimitive curve classes on \(S\). Passing from K3 to \(K3\times E\)
requires the reduced \(E\)-translation quotient and the OP/DT scalar
normalisation above.

\subsection*{Reduced DT and the ring of dual numbers}

\emph{Primary source.} G.~Oberdieck and J.~Shen,
\emph{Reduced Donaldson-Thomas invariants and the ring of dual
numbers}, \emph{Proc.\ Lond.\ Math.\ Soc.} \textbf{118} (2019),
191-220, arXiv:1612.03102.

\emph{Theorem 1 (Oberdieck-Shen).}
Let \(X=S\times E\), with \(S\) a nonsingular
projective K3 surface. Oberdieck and Shen compute all \(E\)-reduced
Donaldson-Thomas invariants in fibre classes \((0,d)\). If
\[
  \sum_{d\ge0}\sum_{n\in\Z} m(d,n)p^n t^d
  =
  -\,\frac{24\,\wp(p,t)}
          {\prod_{m\ge1}(1-t^m)^{24}},
\]
then their product formula gives
\[
  \exp\!\left(
    \sum_{n\ge1}
    \DT^{X,E\text{-red}}_{n,(0,d)}(-p)^n
  \right)
  =
  \prod_{\ell\ge1}
  (1-p^\ell)^{-m(d,\ell)}
\]
for every \(d\ge0\), with the corresponding logarithmic closed formula
obtained by summing over divisors of \((n,d)\).

\emph{Conjecture 1 recalled by Oberdieck-Shen.}
Oberdieck-Pandharipande's multiple-cover
formula for \(K3\times E\) predicts, for \((\beta,n)\neq(0,0)\),
\[
  (-1)^n\DTred_{n,(\beta,d)}
  =
  \sum_{k\mid(n,\beta)}
  \frac{1}{k}\,
  m\!\left(\frac{(\beta/k)^2}{2}+1,\ d,\ \frac{n}{k}\right),
\]
where \(m(h,d,n)\) is read from \(\chiTen^{-1}\). For primitive
\(\beta^2=2h-2\) this gives the Igusa coefficient \(m(h,d,n)\); for
imprimitive \(\beta\) it is a genuine multiple-cover formula.

\emph{Boundary.} Oberdieck-Shen prove the degenerate fibre-class
sector \(\beta=0\) and special imprimitive checks through
\(A\)-equivariant motivic Hall algebras and integration to the ring of
dual numbers. The full imprimitive \(K3\times E\) formula remains the
Oberdieck-Pandharipande multiple-cover conjecture.

\subsection*{Native chiral and Hall boundaries}

For a Calabi-Yau input of dimension \(d\ge3\), the native
CY-to-chiral output is an \(\Eone\)-chiral algebra. An \(\Etwo\) algebra
enters after a centre, Drinfeld double, or module-category layer:
\[
  A_X\in\mathrm{Chir}_{\Eone},
  \qquad
  Z(\mathrm{Rep}(A_X))\in\mathrm{Chir}_{\Etwo}.
\]
For the local affine model, the critical cohomological Hall algebra is
the positive half
\[
  \CoHA(\C^3)=\Yplus(\widehat{\mathfrak{gl}}_1).
\]
The full \(\Woneinfty\) object is obtained after double, centre, completion,
and evaluation data have been supplied. On the \(K3\times E\) branch,
the Hall-Drinfeld/\(\BKM\) denominator-recognition branch is distinct
from the Mukai Yangian self-mirror branch.

\subsection*{Theorem and conjecture loci}

\begin{itemize}
\item \(N=1\), primitive K3 class: Oberdieck-Pixton Theorem~1 proves
  the primitive reduced GW Igusa identity; the OP scalar convention is
  \(\ZOP=-(\PhiTenOP)^{-1}=-4096\,\DeltaFive^{-2}\).
\item Primitive \(\BKM\) denominator: Gritsenko-Nikulin-Borcherds gives
  \(\Dfive=64^{-1}\DeltaFive\) and
  \(\kBKM(\DeltaFive)=c_1(0)/2=5\).
\item Primitive compact CHL denominator rows:
  \((c_N(0))_{N=(1,2,3,4,6)}=(10,8,6,4,2)\) and
  \((\kBKM(\Phi_N))=(5,4,3,2,1)\).
\item \(N\ge2\) CHL primitive scalar: Bryan-Oberdieck Conjecture A
  predicts \(-\widetilde\Phi_N^{-1}\); their Theorem~0.1 proves the
  first \(t\)- and \(q\)-coefficients.
\item CHL dyon sector: Chattopadhyaya-David construct and test the
  Mathieu-twined Siegel modular forms as \(1/4\)-BPS index
  denominators.
\item K3 imprimitive curve counts: Pandharipande-Thomas prove KKV on
  K3 for all curve classes.
\item \(K3\times E\) imprimitive reduced DT: Oberdieck-Pandharipande's
  multiple-cover formula remains conjectural beyond the sectors proved
  by Oberdieck-Shen.
\end{itemize}
